\chapter{平均词、返回矩阵与局部化}\label{sec:arithmetic}

\section{为什么只剩两个参数}
三块核心为
\[
K_{p,r}(a,z)=\bigl(a^p,(rz-a)^p,(-pz)^r\bigr).
\]
总和自动为零。若按照固定计数选择位置，每个新的块值都是 $a,z$ 的有理线性组合。
当若干次平均后重新得到相同的块重数，我们便可以用一个二阶矩阵记录新参数。
这就是返回矩阵的来源。

例如，固定第三块，将前两个 $p$ 块分别按 $p-k,k$ 的互补比例平均，
新第一块值是
\[
a'=\frac{(p-k)a+k(rz-a)}p
   =\frac{p-2k}{p}a+\frac{kr}{p}z,
\qquad z'=z.
\]
它在 $(a,z)$ 坐标中的真实矩阵为
\[
\begin{pmatrix}(p-2k)/p&kr/p\\0&1\end{pmatrix}.
\]
每组取多少份值、是否超出原块的数量，都必须在列出矩阵之前确定。
这种具体选择称为一个平均词：它只是有限多个平均操作按顺序组成的词。

\section{标签、共同尺度与射影返回}

\begin{definition}\label{def:return}
设 $E_\pi:\Q^2\to\Q^n$ 按分组 $\pi$ 将两个参数展开为等值块配置。若平均词 $W$ 满足
\begin{equation}\label{eq:return}
WE_\pi=\lambda E_{\pi'}M,\qquad \lambda\in\Q^\times,
\end{equation}
且 $\pi'$ 与 $\pi$ 具有相同块重数，称 $M$ 是一个可实现返回。构造须对每个同形状分组适用。可逆矩阵的可实现性按非零共同标量识别。
\end{definition}

定义中的 $E_\pi$ 是一个展开映射：例如在某一组标记下，它把 $(a,z)^T$
展开为上式的 $n$ 个数。下标 $\pi$ 记录哪几个原位置属于哪个块。
平均结束后角色可能改变，故输出写成 $E_{\pi'}$。
只要新的块重数相同，就能把下一条构造按新的角色重放。

等式中的共同标量 $\lambda$ 作用于\emph{全部}位置。
由于任意平均词与共同缩放交换，归零问题只依赖参数列的直线方向。
这就是这里“射影”的含义。我们不把乘 $\lambda^{-1}$ 加入允许操作；
实际输出保留这个尺度，后面的归零词在任意非零共同尺度上都同样有效。

\begin{example}[两次正向返回提供逆]
若 $J=\left(\begin{smallmatrix}0&d\\-1&0\end{smallmatrix}\right)$，$d\ne0$，则 $J^2=-dI$。
假如 $J$ 有真实平均词，重复这个词就使全部参数共同乘 $-d$（另计实际词自身的共同尺度）。
所以 $J^{-1}=(-d)^{-1}J$ 在射影意义上由同一条正向词实现。
这不表示一次平均在整个 $\Q^n$ 上可逆；可逆性只发生在带相同块形状的参数空间上。
\end{example}

各返回的分组将由非负整数计数给出。式\eqref{eq:return}允许组合返回；如果 $M^j=cI$、$c\ne0$，其 $j-1$ 次返回实现 $M^{-1}$。


\section{剪切矩阵与可用分母}
上剪切把 $(x,y)$ 送到 $(x+ty,y)$；下剪切把它送到 $(x,y+tx)$。
它们的优点是参数可直接相加：
\[
U(s)U(t)=U(s+t),\qquad U(t)^{-1}=U(-t),
\]
下剪切同理。这里的“根子群”就是固定一种方向、让剪切参数遍历某个加法群所得的矩阵群。

构造中常能先实现一个整数剪切，然后用一个可逆倍率把参数不断除以某些整数。
将这些参数写成一个环比逐个列分母更方便。
\begin{definition}[整数的有限局部化]
对整数 $D>1$，
\[
\Z[1/D]=\{a/D^k:a\in\Z,\ k\ge0\}.
\]
等价地，约分后的分母只含 $D$ 的素因子。
其单位是 $\pm\prod_{\ell\mid D}\ell^{e_\ell}$，其中指数 $e_\ell$ 可正可负。
\end{definition}
例如 $\Z[1/6]$ 包含 $1/2,1/3,1/12$，不包含 $1/5$。
数 $5$ 属于这个环，但不是单位。这个区别会影响能否把某个剩余尺度提升成真实可用的共同单位。

若 $\gcd(q,D)=1$，约化映射把 $a/D^k$ 送到 $a(D^k)^{-1}\pmod q$，
从而 $R/qR$ 与 $\Z/q\Z$ 同构。
若 $q$ 含 $D$ 的素因子，先把这些单位因子删除即可。

记 $U(t)=\left(\begin{smallmatrix}1&t\\0&1\end{smallmatrix}\right)$、$L(t)=\left(\begin{smallmatrix}1&0\\t&1\end{smallmatrix}\right)$，以及
\[
R_D=\Z[1/D],\qquad
\Gamma(q,R_D)=\ker(\SL_2(R_D)\to\SL_2(R_D/qR_D)).
\]
以下每次扩张 $R_D$ 均由引理\ref{lem:localize}给出；$D$ 的素因子不与相应配置的危险素数相交。


\section{一个倍率怎样产生完整的参数环}
声称“有 $U(hR)$”意味着对\emph{每个} $t\in hR$ 都有有限的真实平均词，
并不只是在有限模数上看到了这些矩阵。
下一个引理从已经可逆的三角返回出发，给出扩大参数环的具体办法。

\begin{lemma}[可逆倍率扩展剪切参数]\label{lem:localize}
假设 $U(hR)$ 全部可实现，且上三角矩阵 $T=\left(\begin{smallmatrix}1&b\\0&\kappa\end{smallmatrix}\right)$ 及其逆可实现。则 $U(hR[\kappa,\kappa^{-1}])$ 全部可实现。若 $\kappa=A/B$ 约分，则新环为 $R[1/(AB)]$。
\end{lemma}
\begin{proof}
使用 $TU(t)T^{-1}=U(t/\kappa)$、$T^{-1}U(t)T=U(\kappa t)$ 和 $U(s)U(t)=U(s+t)$。
新环中每个元素是有限个 $r_i\kappa^i$ 的和，其中 $i$ 允许为负数。
先用共轭实现每一项的剪切，再相乘把参数相加，得到所需返回。
全部求和都是有限的，因此最终仍是一条有限平均词。
至于新环中的分母为什么正好来自 $AB$，一方面 $A/B$、$B/A$ 显然都属于 $R[1/(AB)]$；
另一方面，裴蜀等式 $uA+vB=1$ 给出 $u\kappa+v=1/B$；对逆倍率同理得到 $1/A$。
\end{proof}

